\subsection{From simulations to observations}
\label{sec:transfer}

A response constraint is evaluated within the same image model that supplies the training data. Applying the trained estimator to SuperBIT observations introduces a further question: how well does that model represent the images being measured? \emph{Distribution shift} occurs when the observations differ from the examples used for training. For this application, relevant differences include galaxy morphology, neighboring light, PSF structure, centering, masking, and noise correlations. Matching a one-dimensional flux distribution cannot establish that their joint distribution is representative. A model can therefore perform well on held-out simulations generated by the same assumptions and still fail on flight data~\cite{mandelbaum_weak_2018,tewes_weak-lensing_2019}.

Testing this transfer requires both known-shear simulations and diagnostics that can be measured on the real images. Known-shear simulations test $m$ and $c$; real images permit tests of PSF residuals, mean-shape trends, and spatial consistency. Comparing estimators on matched objects helps separate a change in measurement from a change in the selected population. Neither agreement with a reference estimator nor a visually convincing mass peak supplies the missing truth labels. The resulting catalog must therefore be assessed through both calibration tests and observable diagnostics before its shapes are used to reconstruct the cluster mass distribution~\cite{saha_lensing_2026-1,mandelbaum_weak_2018}.
